\section{Практическая часть}

\subsection{Цель работы}

Целью работы является изучение различных режимов возбуждения лампового генератора и установление зависимости амплитуды и частоты автоколебаний от параметров системы.

\subsection{Экспериментальная установка}

Схема лабораторной установки приведена на \cref{fig:pr:1}.

\begin{figure}[H]
	\centering
	\includegraphics[width=0.7\linewidth]{pr_1}
	\caption{Схема лабораторной установки}
	\label{fig:pr:1}
\end{figure}

Для исследования фазовой плоскости на осциллограф подаются сигналы, соответствующие напряжениям между точками 1--3 и 2--3. Коэффициент взаимоиндукции $M$ используется как параметр обратной связи, а напряжение смещения сетки позволяет менять режим работы генератора.

\subsection{Задания, выполняемые при подготовке}

\begin{enumerate}
	\item Познакомиться с методом медленно меняющихся амплитуд [1, 2].
	\item Составить исходные и укороченные уравнения исследуемой системы.
	\item Найти состояния равновесия исходных и укороченных уравнений, исследовать их устойчивость и определить их типы.
	\item Теоретически построить бифуркационные диаграммы и фазовые портреты для различных режимов работы.
	\item Познакомиться с экспериментальной установкой \cref{fig:pr:1}.
\end{enumerate}

\subsection{Задания, выполняемые в лаборатории}

\subsubsection{Мягкий режим}

Подбором напряжения смещения необходимо получить режим, в котором увеличение коэффициента взаимоиндукции сопровождается плавным ростом амплитуды установившихся автоколебаний, начиная с нуля.

Снять бифуркационную диаграмму зависимости амплитуды автоколебаний от коэффициента взаимоиндукции $M$.

\subsubsection{Сложно-жёсткий режим}

Подобрать напряжение смещения, соответствующее сложно-жёсткому типу генератора, снять бифуркационную диаграмму.

\subsection{Протокол измерений}

\begin{table}[H]
	\centering
	\caption{Мягкий режим}
	\label{tab:pr:1}
	\begin{tabular}{cc}
		\toprule
		Коэффициент взаимоиндукции $M$ & Амплитуда $U$, В \\
		\midrule
		10 & 15 \\
		20 & 28 \\
		30 & 35 \\
		40 & 38 \\
		50 & 40 \\
		65 & 41 \\
		90 & 43 \\
		110 & 45 \\
		130 & 46 \\
		\bottomrule
	\end{tabular}
\end{table}

\begin{table}[H]
	\centering
	\caption{Сложно-жёсткий режим}
	\label{tab:pr:2}
	\begin{tabular}{cc}
		\toprule
		Коэффициент взаимоиндукции $M$ & Амплитуда $U$, В \\
		\midrule
		10 & 0 \\
		57 & 39 \\
		70 & 41 \\
		90 & 43 \\
		110 & 44 \\
		130 & 45 \\
		120 & 43 \\
		110 & 42 \\
		90 & 41 \\
		70 & 41 \\
		60 & 39 \\
		50 & 35 \\
		40 & 27 \\
		35 & 25 \\
		32 & 20 \\
		30 & 0 \\
		\bottomrule
	\end{tabular}
\end{table}

\subsection{Построение практических графиков}

\textbf{Задание 1.} Снятие бифуркационной диаграммы при мягком типе генератора (\cref{fig:pr:2a,fig:pr:2b}).

\begin{figure}[H]
	\centering
	\begin{minipage}{0.45\textwidth}
		\centering
		\includegraphics[width=\linewidth]{pr_2}
		\subcaption{Экспериментальные данные}
		\label{fig:pr:2a}
	\end{minipage}
	\begin{minipage}{0.45\textwidth}
		\centering
		\includegraphics[width=\linewidth]{pr_3}
		\subcaption{Аппроксимация}
		\label{fig:pr:2b}
	\end{minipage}
	\caption{}
	\label{fig:pr:2}
\end{figure}

\textbf{Задание 2.} Снятие бифуркационной диаграммы при сложно-жёстком типе генератора (\cref{fig:pr:3a,fig:pr:3b}).

\begin{figure}[H]
	\centering
	\begin{minipage}{0.45\textwidth}
		\centering
		\includegraphics[width=\linewidth]{pr_4}
		\subcaption{Экспериментальные данные}
		\label{fig:pr:3a}
	\end{minipage}
	\begin{minipage}{0.45\textwidth}
		\centering
		\includegraphics[width=\linewidth]{pr_5}
		\subcaption{Аппроксимация}
		\label{fig:pr:3b}
	\end{minipage}
	\caption{}
	\label{fig:pr:3}
\end{figure}
